\chapter{Surgical Aortic Valve Replacement}

\section*{Overview}
Surgical aortic valve replacement (SAVR) replaces a diseased aortic valve with a mechanical or tissue prosthesis. It is distinct from transcatheter aortic valve replacement (TAVR).

\section*{Alternative Names}
Open aortic valve replacement, SAVR, aortic valve surgery

\section*{Why This Test or Procedure Is Ordered}
SAVR is used for severe symptomatic aortic stenosis, severe aortic regurgitation, or a diseased valve requiring replacement when surgical treatment is appropriate.

\section*{How This Test or Procedure Is Performed}
Under general anesthesia, the chest is opened through a sternotomy or another surgical approach. A heart-lung bypass machine maintains circulation while the native valve is removed and the prosthesis is implanted.

\section*{Preparation, Risks, and Recovery}
Evaluation includes echocardiography and other cardiac testing as indicated. Risks include bleeding, infection, stroke, rhythm problems, kidney injury, blood clots, valve dysfunction, and death. Recovery begins in intensive care and continues with mobilization and cardiac rehabilitation. Mechanical valves generally require lifelong anticoagulation.

\section*{Follow-up Tests and Procedures}
Postoperative echocardiography and regular cardiology review are required. Mechanical-valve recipients need anticoagulation monitoring.

\section*{References}
\begin{enumerate}
  \item MedlinePlus. Aortic Valve Surgery---Open. U.S. National Library of Medicine; 2025.
  \item NHS England. TAVI and Surgical Aortic Valve Replacement Position Statement; 2023.
\end{enumerate}

\clearpage
